\documentclass[10pt,a4paper]{amsart}
\usepackage[margin=19mm]{geometry}
\usepackage{amsmath,amssymb,mathtools}
\usepackage[scr=rsfs,cal=euler]{mathalfa}
\usepackage{xfrac}
\usepackage[unicode,hidelinks,bookmarks=false]{hyperref}
\newcommand{\ie}{i.e.\ }
\newcommand{\cf}{cf.\ }
\newcommand{\resp}{respectively\ }
\newcommand{\Subsection}{\subsection}
\providecommand{\nobreakdash}{\nobreak\hbox{-}}
\setlength{\emergencystretch}{2em}
\multlinegap0pt
\usepackage{amssymb}
\newtheorem{lemma}{Lemma}
\title{\bf {A necessary condition for solvability by radicals}}
\author{Askold Khovanskii}
\date{Source: arXiv:2604.08642v1 (CC BY 4.0)}
\begin{document}
\maketitle

\noindent\emph{Source and licence.} \bf {A necessary condition for solvability by radicals}. Version arXiv:2604.08642v1 (CC BY 4.0), distributed by arXiv under the Creative Commons Attribution 4.0 International licence, http://creativecommons.org/licenses/by/4.0/, as recorded in the arXiv licence metadata for this exact version and shown on the abstract page https://arxiv.org/abs/2604.08642v1. Full licence notice: \texttt{licenses/arxiv-2604.08642-CC-BY-4.0.txt}.\par\smallskip

\noindent\textit{Editorial scope.} Counted unit: the article's lemma rroot (source0.tex lines 240--241). The article's complete original proof of it is delivered with it (source0.tex lines 243--250, one \texttt{proof} environment of 8 lines). No mathematics is written, repaired or re-derived: the statement and the proof are the article's own lines, in the article's own notation, and no retained line is substituted or re-flowed.  No further source-local context is needed: every cross-reference of every retained line resolves inside the two delivered spans.  Nothing was omitted from any retained span, so the package carries no omission record.  No retained line contains a citation, so the wrapper carries no bibliography and no bibliography entry.  No retained line contains a figure environment, an image-inclusion command or drawing code, so no image is delivered and the package declares no asset dependency.  Editorial formatting only, in addition: an AMS \texttt{amsart} preamble that restates the article's own declarations and macros used by the retained lines, namely the environments \texttt{lemma} on separate counters, together with the standard packages those lines need.  The retained environments print in this document as Lemma 1 in order of appearance, whereas the article prints the same items with the numbering of its own full document.  The complete native source of the article as distributed in the arXiv e-print for this exact version is bundled unchanged as \texttt{sources/198165/source0.tex}.
\begin{lemma}\label{rroot} Let $E$ be a splitting field over a ground field $K$ of a polynomial $x^n-a$, where $n>1$  is a natural number and $a\in K$. Assume that $K$ contains all roots of unity of degree $n$.  Then the Galois group $G(E,K)$ is isomorphic to a subgroup of  the cyclic additive group $\Bbb Z_n$. In particular, the group $G(E,K)$ is commutative.
\end{lemma}

\begin{proof} Let $x_0$ be a root of the equation $x^n-a=0$. Each  root of this equation can be represented as $x_0 \omega^k$, where $\omega$ is a generator of the multiplicative group of  roots of unity of degree $n$, and the integer $k$ is defined up to addition of a number divisible by $n$.
Thus, one can identify the set of roots of the polynomial $x^n-a$ with the set of integral numbers modulo $n$.

If an element $g\in G(E,K)$ sends
the root $x_0$ to a root $x_0\omega^m$, then it sends a root $x_0 \omega^k$ to the root $x_0\omega^k \omega^m=x_0\omega^{m+k}$ for every integer $k$.  Indeed, by the assumption,    $\omega\in K$, so  we have  $g(\omega)=\omega$ and $g(\omega^k)=\omega^k$.   Thus, under the identification of the set of roots of the polynomial $x^n-a$ with integers modulo $n$, the  action of the element $g$ is  reduced to the addition of the element  $(m\mod n)$ to elements  $k\in \Bbb Z_n$. We obtain a map from the Galois group $G(E,K)$ to the group $\Bbb Z_n$ which sends an element $g$ to an element $m$. One can see that this map is a group homomorphism. This homomorphism has no kernel since the field $E$ is generated over the field $K$ by roots of the equation $x^n-a=0$.
Thus, the homomorphism embeds the Galois group $G(E,K)$ to the group $\Bbb Z_n$.

\end{proof}

\end{document}
